% 问题二：总体路线

本题是一个\textbf{约束满足 + 多目标优化}问题。由于搜索空间大（每个计划有多种调整方式和撤销选项），单一算法难以保证全局最优，因此采用\textbf{4 种方法求解并对比}的策略：

\begin{enumerate}
    \item \textbf{贪心 + 优先级}：按 A > B > C 优先级排序，逐个计划尝试保留、调整、撤销，每步选择最小幅度的可行动作。速度快，作为基准方法。
    \item \textbf{整数线性规划（ILP）}：将问题建模为 0-1 整数规划，以加权撤销数（B 类权重高于 C 类，A 类固定保留）为目标函数，用 CBC 求解器求解。理论上可得全局最优，但规模较大时可能超时。
    \item \textbf{贪心构造（对照实现）}：采用与方案 1 相同的贪心构造逻辑、按幅度从小到大递增选择动作的独立实现，作为方案 1 的交叉验证，用于确认贪心结果的稳健性。
    \item \textbf{模拟退火（SA）}：以贪心解为初始解，通过随机扰动（随机更换 1\textasciitilde3 个计划的动作）探索解空间，Metropolis 准则接受劣解以跳出局部最优。采用多种子策略（$\geq 5$ 个种子），收尾用贪心消毒保证 0 冲突。
\end{enumerate}

最终选择撤销数最少的方案作为最优解。
